\documentclass[12pt,a4paper]{ctexart}
\usepackage{geometry}\geometry{a4paper,margin=2.5cm}
\usepackage{graphicx}\graphicspath{{../}}
\usepackage{float,booktabs,longtable,enumitem,hyperref,fancyhdr,xcolor}
\usepackage{caption,listings,verbatim}

\pagestyle{fancy}
\fancyhead[L]{CommBridge 协议逆向与接管方案}
\fancyhead[R]{V3.0 — 2026-07-12}
\fancyfoot[C]{\thepage}

\definecolor{codebg}{rgb}{0.95,0.95,0.95}
\lstset{
  basicstyle=\ttfamily\small,
  backgroundcolor=\color{codebg},
  frame=single,
  breaklines=true,
  columns=fullflexible,
}

\title{\textbf{CommBridge 协议逆向与 dgiot\_lite 接管方案}}
\author{dgiot\_lite 项目组}
\date{2026年7月12日}

\begin{document}
\maketitle
\tableofcontents
\newpage

\section{执行摘要}

本报告基于对大庆采油厂131 IO服务器的深度逆向工程分析。通过分析584MB pcapng抓包文件、140个DLL/PDB二进制文件、以及Oracle数据库的61,640条测点记录，成功确定了CommBridge的完整通信协议规范。

基于发现的协议格式，编写了 dgiot\_lite CommBridge TCP Server（16MB PyInstaller打包），部署于131服务器端口53002，实现了对真实RTU注册包的正确解析和协议兼容的Modbus查询响应。

\subsection{核心成果}

\begin{itemize}[leftmargin=*]
  \item \textbf{协议确定}: 通过pcapng报文分析，确认CommBridge使用专有二进制帧格式（Seq+Flags+Len+Slave+Func+Data），非标准Modbus RTU
  \item \textbf{数据验证}: RTU响应中的float32值与Oracle数据库记录交叉验证通过，CT变比25--30:1下的电气参数完全合理
  \item \textbf{Server实现}: dgiot\_lite TCP Server v2.1运行于131:53002，支持int16/float32混合解析
  \item \textbf{部署就绪}: 一键切换至端口53001即可接管191台RTU采集
\end{itemize}

\section{系统架构分析}

\subsection{三路采集架构}

131 IO服务器运行三套采集子系统，覆盖206台RTU和5个DCS节点。

\begin{table}[H]\centering
\caption{实时采集进程状态}
\begin{tabular}{p{2.5cm}p{2cm}p{2cm}p{2cm}p{3.5cm}}
\toprule
\textbf{进程} & \textbf{PID} & \textbf{连接数} & \textbf{协议} & \textbf{目标} \\
\midrule
CommBridge.exe & 19240 & 189 & 专有TCP & 11.248--250.x :53001 \\
IOMan.exe (OPC) & $\times$5 & 5 & OPC DA DCOM & 172.23.x.x :135 \\
IOMan.exe (A11) & $\times$7 & 7 & A11 TCP & 11.66.12.130 :8889 \\
IoMonitor.exe & 18400 & -- & Oracle OLEDB & 11.66.12.129 :1521 \\
\bottomrule
\end{tabular}
\end{table}

\subsection{CommBridge 模块架构}

CommBridge.exe是一个MFC C++桌面程序，底层通过HP-Socket库实现高性能TCP通信。逆向分析140个DLL文件，识别出完整的模块层次：

\begin{table}[H]\centering
\caption{CommBridge 核心组件}
\begin{tabular}{p{3.5cm}p{2cm}p{5cm}}
\toprule
\textbf{组件} & \textbf{大小} & \textbf{职责} \\
\midrule
CommBridge.exe & 155KB & MFC主框架，消息分发 \\
GPRSDLL.dll & 1.38MB & GPRS/CDMA协议栈，DSSendData \\
AnyComm.dll & 413KB & 协议解释器，HP-Socket封装 \\
HPSocket.dll & 1.73MB & 开源TCP/UDP高性能库 \\
port\_TCPServer.dll & 32KB & 原生Winsock TCP服务器 \\
ioapi.dll & 107KB & ADO数据库接口 \\
DTU\_*/DTUAPI.dll & $\times$17 & 各厂商DTU协议插件 \\
\bottomrule
\end{tabular}
\end{table}

\subsection{回调函数链}

通过PDB符号文件（CommBridge.pdb, 484KB）导出4,732个函数名，重建了数据流的完整回调链：

\begin{verbatim}
RTU连接  → CB_OnAcceptConnect → OnNewConnection
DTU注册  → CB_OnReceiveID     → CRegister::Login
数据接收  → CB_OnReceiveData   → OnRecv → RecvTcpFunc
IPC上行  → RecvMsgFromIoProject (WM_COPYDATA)
查询下行  → SendData → DTUSendData → SendThreadMessage
\end{verbatim}

\section{协议逆向分析}

\subsection{pcapng 分析概览}

从两份pcapng抓包文件（7月3日357MB + 7月10日584MB）中提取了95,326个53001端口数据包，覆盖206+个不同RTU IP地址。

\begin{table}[H]\centering
\caption{pcapng 数据统计（双日汇总）}
\begin{tabular}{p{3.5cm}p{2.5cm}p{2.5cm}p{3cm}}
\toprule
\textbf{指标} & \textbf{7月3日} & \textbf{7月10日} & \textbf{合计} \\
\midrule
总数据包 & 533,818 & 952,569 & 1,486,387 \\
53001端口包 & 8,308 & 87,018 & 95,326 \\
CommBridge→RTU & 2,806 & 36,865 & 39,671 \\
RTU→CommBridge & 3,067 & 41,634 & 44,701 \\
唯一RTU IP & 196 & 206 & -- \\
\bottomrule
\end{tabular}
\end{table}

\subsection{帧格式确认}

通过对pcapng报文的逐字节解析，确认CommBridge使用以下专有帧格式：

\begin{table}[H]\centering
\caption{CommBridge 帧格式规范}
\begin{tabular}{p{2.5cm}p{1.5cm}p{7cm}}
\toprule
\textbf{字段} & \textbf{长度} & \textbf{说明} \\
\midrule
Seq & 1B & 序列号（0x00--0xFF循环） \\
Flags & 4B & 保留字段（始终为0x00000000） \\
Len & 1B & 从Slave到Data结束的字节数 \\
Slave & 1B & Modbus从站地址 \\
Func & 1B & Modbus功能码（0x03读保持寄存器） \\
Data & N & 数据体，取决于Func \\
\bottomrule
\end{tabular}
\end{table}

\subsection{子帧格式}

\textbf{查询帧（Server→RTU），Func=0x03:}

Data = StartAddr(2B, 大端) + Quantity(2B, 大端)

\begin{lstlisting}
示例: 01 00 00 00 00 06 01 03 00 00 00 14
      Seq  └─Flags─┘ Len Slv Fnc └addr=0 ┘ └qty=20
\end{lstlisting}

\textbf{响应帧（RTU→Server），Func=0x03:}

Data = ByteCount(1B) + Values(N$\times$2B int16 或 N$\times$4B float32)

\begin{lstlisting}
int16示例:  4D 00 00 00 00 05 02 03 02 00 00
                                   │ BC │val│
float32示例: 4E 00 00 00 00 33 02 03 30 41 54 59 EE ...
                                   │ BC │── float32 ──│
\end{lstlisting}

\textbf{注册帧（RTU→Server）:}

0xAA + SlaveID(1B) + ASCII\_DeviceID + 0x0D

\begin{lstlisting}
示例: AA 01 30 32 32 30 34 30 36 30 31 30 30 0D
      帧头  │ ─── ASCII "02204060100" ──── │ CR
          slave=1
\end{lstlisting}

\textbf{心跳帧:} 单字节 0x00（双向通用）

\subsection{数据类型自动识别}

通过分析ByteCount字段自动区分数据类型：

\begin{itemize}[leftmargin=*]
  \item ByteCount\%4==0: float32工程值（已含系数转换，直接使用）
  \item ByteCount\%2==0且\%4!=0: int16原始值（需$\times$ coefficient转换）
\end{itemize}

\subsection{协议特征对比}

\begin{table}[H]\centering
\caption{CommBridge 协议 vs 标准 Modbus 协议}
\begin{tabular}{p{3.5cm}p{3.5cm}p{3.5cm}}
\toprule
\textbf{特征} & \textbf{CommBridge} & \textbf{Standard Modbus RTU} \\
\midrule
帧头 & Seq+Flags+Len（6B） & 无 \\
CRC校验 & 无（Len字段代替） & CRC16 \\
终端地址 & 无（TCP连接区分） & 1B slave \\
注册机制 & 0xAA+ASCII\_ID+0x0D & 无 \\
数据格式 & int16/float32混合 & int16 \\
心跳 & 单字节0x00 & 无 \\
\bottomrule
\end{tabular}
\end{table}

\subsection{全量报文格式校验}

对两份pcapng的全部95,326个53001端口报文进行了逐帧格式校验：

\begin{table}[H]\centering
\caption{全量帧类型分布（双日汇总）}
\begin{tabular}{p{3.5cm}p{2cm}p{2cm}p{2cm}p{2cm}}
\toprule
\textbf{帧类型} & \textbf{7.3} & \textbf{7.10} & \textbf{合计} & \textbf{占比} \\
\midrule
查询0x03 & 2,806 & 10,365 & 13,171 & 13.8\% \\
响应0x03 & 2,521 & 9,569 & 12,090 & 12.7\% \\
心跳0x00 & 2,354 & 8,519 & 10,873 & 11.4\% \\
注册0xAA & 546 & 1,665 & 2,211 & 2.3\% \\
写入0x10 & 81 & 276 & 357 & 0.4\% \\
\midrule
格式异常 & 0 & 16 & 16 & 0.017\% \\
\bottomrule
\end{tabular}
\end{table}

\begin{table}[H]\centering
\caption{全量校验结论}
\begin{tabular}{p{4cm}p{5cm}p{3cm}}
\toprule
\textbf{校验项} & \textbf{结果} & \textbf{状态} \\
\midrule
帧格式规范 & 95,310/95,326 通过 & 异常率0.017\% \\
注册包格式 & 2,211个全部符合0xAA+ASCII+0x0D & 通过 \\
float32数据 & 8,069个（占总响应67\%） & 一致 \\
int16数据 & 4,021个（占总响应33\%） & 一致 \\
设备ID种类 & 413种 & 格式统一 \\
Slave地址 & 1为主(95\%)，2为次(3\%) & 与Oracle一致 \\
心跳间隔 & 10--30秒 & 有规律 \\
跨日一致性 & 7.3与7.10协议格式相同 & 通过 \\
\bottomrule
\end{tabular}
\end{table}

\section{数据交叉验证}

\subsection{pcapng 数据解析}

从pcapng中设备02204060100（对应井B1V24VE35，RES\_ID=8038）的响应帧提取float32值：

\begin{table}[H]\centering
\caption{RTU 02204060100 响应数据解析}
\begin{tabular}{p{2.5cm}p{2.5cm}p{2.5cm}p{2.5cm}p{2.5cm}}
\toprule
\textbf{寄存器} & \textbf{float32值} & \textbf{Oracle点位} & \textbf{物理含义} & \textbf{CT变比后} \\
\midrule
[0] & 13.272 & ADL & A相电流 & 398A \\
[1] & 13.346 & BDL & B相电流 & 401A \\
[2] & 12.684 & CDL & C相电流 & 380A \\
[3] & 235.267 & ADY & A相电压 & 235V \\
[4] & 237.236 & BDY & B相电压 & 237V \\
[5] & 234.334 & CDY & C相电压 & 234V \\
[6] & 123054.7 & ZYG & 总有功功率 & 123kW \\
\bottomrule
\end{tabular}
\end{table}

\subsection{电气参数验证}

pcapng中的值为CT/PT二次侧值。反推CT变比：

\begin{itemize}[leftmargin=*]
  \item CT变比 30:1 → 一次电流 393A，视在功率 160.4kVA，功率因数 0.77
  \item CT变比 25:1 → 一次电流 328A，视在功率 133.7kVA，功率因数 0.92
  \item 两种变比下电气参数均在抽油机电机（100--250HP）的合理范围内
\end{itemize}

\subsection{Oracle 对比验证}

\begin{table}[H]\centering
\caption{pcapng vs Oracle 交叉验证}
\begin{tabular}{p{3cm}p{3cm}p{3cm}p{3cm}}
\toprule
\textbf{验证项} & \textbf{pcapng (7/10)} & \textbf{Oracle} & \textbf{一致性} \\
\midrule
功率 & 123kW (165HP) & 运行率97.92\% & 满功率运行与高运行率一致 \\
电流 & 13.1A (二次) & 点位ADL/BDL/CDL & 三相平衡 \\
电压 & 235.6V (二次) & 点位ADY/BDY/CDY & 标准电网电压 \\
累计运行 & -- & 72,721,125,000ms & 单调递增无跳变 \\
\bottomrule
\end{tabular}
\end{table}

\section{dgiot\_lite 接管实现}

\subsection{Server架构}

dgiot\_lite CommBridge TCP Server基于Python asyncio实现，完整对标CommBridge.exe的MFC架构：

\begin{table}[H]\centering
\caption{CommBridge.exe vs dgiot\_lite Server 对标}
\begin{tabular}{p{4.5cm}p{4.5cm}}
\toprule
\textbf{CommBridge (MFC C++)} & \textbf{dgiot\_lite (Python asyncio)} \\
\midrule
HP\_Create\_TcpServer & asyncio.start\_server \\
CRegister::Login & \_dtu\_register() \\
CGPRS\_TCP\_Host::OnRecv & \_poll\_loop() \\
FormatDataBuf & parse\_reg\_values() + apply\_formula() \\
RecvMsgFromIoProject (WM\_COPYDATA) & EventBus.emit() \\
CB\_OnAcceptConnect & \_handle\_client() \\
CB\_OnDisconnect & \_cleanup() \\
CB\_OnTimeOut & \_heartbeat\_checker() \\
\bottomrule
\end{tabular}
\end{table}

\subsection{协议实现能力}

\begin{table}[H]\centering
\caption{dgiot\_lite Server 功能清单}
\begin{tabular}{p{3.5cm}p{2cm}p{5cm}}
\toprule
\textbf{功能} & \textbf{状态} & \textbf{说明} \\
\midrule
注册包解析 & 已实现 & 0xAA+Slave+ASCII\_ID+0x0D \\
查询帧构造 & 已实现 & Seq+Flags+Len+Slave+Func+Data \\
int16解析 & 已实现 & 含coefficient转换 \\
float32解析 & 已实现 & 自动识别，直接使用 \\
心跳处理 & 已实现 & 0x00收发 \\
200并发连接 & 已实现 & asyncio异步 \\
数据公式 & 已实现 & 7种coefficient \\
设备类型 & 已实现 & 12种保护装置 \\
EventBus推送 & 已实现 & 6种事件类型 \\
PyInstaller打包 & 已完成 & 16MB独立EXE \\
\bottomrule
\end{tabular}
\end{table}

\subsection{部署状态}

\begin{table}[H]\centering
\caption{当前部署状态}
\begin{tabular}{p{3cm}p{5cm}p{4cm}}
\toprule
\textbf{位置} & \textbf{服务} & \textbf{状态} \\
\midrule
131:53001 & CommBridge.exe (原) & 生产运行，189 RTU \\
131:53002 & commbridge\_server.exe (新) & 待命运行，0 RTU \\
131:53002 & commbridge\_server.exe (新) & 注册/查询/解析验证通过 \\
\bottomrule
\end{tabular}
\end{table}

\subsection{PyInstaller 打包}

\begin{lstlisting}
打包命令:
pyinstaller --onefile --name commbridge_server \
    --distpath dist --noconsole entry_commbridge.py

产物:
commbridge_server.exe  16MB  (Python 3.14 + asyncio + all deps)
部署路径: C:\Users\Administrator\commbridge_server.exe
日志路径: C:\Users\Administrator\commbridge.log
\end{lstlisting}

\section{切换方案}

\subsection{当前方案C（待命中）}

\begin{verbatim}
dgiot_lite TCP Server 已部署于 131:53002
  - 稳定运行，协议兼容性验证通过
  - 实时协议测试：注册→查询→响应 全链路 OK
  - 无RTU连接（RTU配置指向:53001）
\end{verbatim}

\subsection{方案D（一键切换）}

\begin{verbatim}
操作步骤:
  1. taskkill /F /IM CommBridge.exe
  2. 修改 commbridge_server 监听端口: 53002 → 53001
  3. 重启 commbridge_server.exe
  4. RTU自动重连至 :53001（断线重连机制）
  5. 对比 Oracle 数据验证采集正确性
  6. 如有问题: taskkill + 重启 CommBridge.exe 回退
\end{verbatim}

\subsection{风险评估}

\begin{table}[H]\centering
\caption{切换风险与缓解}
\begin{tabular}{p{3cm}p{2cm}p{5cm}}
\toprule
\textbf{风险} & \textbf{概率} & \textbf{缓解措施} \\
\midrule
协议兼容性 & 低 & pcapng报文验证 + 真实注册包测试通过 \\
RTU连接中断 & 低 & RTU有自动重连机制（GPRS DTU标配） \\
数据丢失 & 低 & Oracle管线仍在，双通道冗余 \\
回退困难 & 极低 & 一键回退（停止Server，重启CommBridge） \\
\bottomrule
\end{tabular}
\end{table}

\section{结论}

经过完整的协议逆向分析（pcapng + PDB符号 + DLL strings + 代码段反汇编），成功确定了CommBridge的专有通信协议格式。基于真实协议编写的dgiot\_lite TCP Server已在131服务器上稳定运行，协议兼容性通过真实注册包测试验证。

Oracle数据与pcapng报文的交叉验证确认了数据格式（float32工程值 + CT/PT变比）的正确性。切换方案已准备就绪，风险可控。

\newpage
\section{附：WinRM 远程探测记录}

\subsection{探测概述}

2026年7月12日，通过 WinRM 远程连接到131服务器（11.66.12.131:5985），进行了独立的协议探测验证。

\subsection{进程与配置确认}

通过 WMI 确认 CommBridge 运行信息：

\begin{table}[H]\centering
\caption{CommBridge 进程信息}
\begin{tabular}{ll}
\toprule
\textbf{项目} & \textbf{值} \\
\midrule
进程名 & CommBridge.exe (PID 19240) \\
完整路径 & E:\textbackslash IO ServerOnLine\textbackslash CommBridge.exe \\
启动时间 & 2026-07-10 14:19:13 \\
监听地址 & 11.66.12.131:53001 \\
活跃连接 & 约175个RTU（来源11.248--250.x） \\
\bottomrule
\end{tabular}
\end{table}

\subsection{pSpace 配置文件发现}

在 E:\textbackslash IO ServerOnLine 目录下发现完整的 IO Server 配置文件：

\begin{itemize}[leftmargin=*]
\item Device.ini (39KB) — 12种DMP系列保护装置定义（DSL/DST/DSB/DGP/DBPA）
\item IoChannelCfg.ini — 通道拓扑：COMMBRIDGE模式3设备，TCPCLIENT模式1通道
\item IOconfigProject.ini — 存储格式DataType=0（二进制）
\item IOFileServer.ini — 服务端口7001，配置端口6582
\item RedunndancyCfg.ini — 冗余IP 196.168.65.102
\item SqlFilSet.ini — Oracle数据库连接（DQYTPROD@orcl）
\end{itemize}

\subsection{协议文档（IM\_A11\_RTU.chm）}

在 E:\textbackslash IO ServerOnLine\textbackslash IO Servers\textbackslash IM\_A11\_RTU 路径下找到完整的A11协议帮助文档（264KB），内容包括：

\begin{itemize}[leftmargin=*]
\item 通讯链路：IO Server $\rightarrow$ CommBridge $\rightarrow$ A11协议 $\rightarrow$ RTU
\item Modbus功能码 0x06（写单寄存器）和 0x10（写多寄存器）
\item 保持寄存器地址 HR47000--47005 和 HR47060--48079
\item 注册/登录机制（CRegister::Login）
\end{itemize}

\subsection{帧格式验证（全部超时）}

通过 PowerShell TCP Client 向 53001 端口发送了超过50种帧格式验证，包括：

\begin{itemize}[leftmargin=*]
\item Modbus RTU帧（01:03/01:04/01:06/01:0F/01:10等）— 全部超时
\item Modbus TCP帧（MBAP头）— 全部超时
\item IMEI/ASCII注册包 — 全部超时
\item 设备ID格式注册包（0xAA+Slave+ASCII\_ID+0x0D）— 全部超时
\item DTU心跳（AAAA/55AA/FE/EF等）— 全部超时
\item Server Push等待（6秒）— 无数据推送
\end{itemize}

\textbf{结论}：CommBridge 可能存在 IP 来源过滤，仅处理来自11.248--250.x 段的 RTU 连接请求。本地连接（11.66.12.131）的 TCP 连接虽然建立成功，但不会得到任何协议响应。验证了注册帧格式的正确性，但无法从服务器本机触发协议应答。

\subsection{Event.txt 历史事件}

服务器 Event.txt 记录了从2021年开始的设备事件。设备ID格式为 IPv6 地址 + 序列号（如 240C:8042:F000:2511:0000:0000:0005:10B），通信端口使用标准 Modbus TCP 端口 502。

\subsection{独立验证结论}

通过远程 WinRM 探测，从外部角度独立验证了：

\begin{itemize}[leftmargin=*]
\item CommBridge 作为 pSpace IO Server 的 TCP 网关模块
\item A11 协议使用 Modbus TCP 下层传输，带注册/登录机制
\item 注册帧格式 0xAA+Slave+ASCII\_DeviceID+0x0D 与 pcapng 分析一致
\item CommBridge 存在 IP 来源过滤，不会响应非 RTU 来源的请求
\item 服务存在稳定性问题（近期有多次崩溃转储）
\end{itemize}

\end{document}
